\chapter{Introduction}
\label{sec:introduction}

The task of deleting data and ensuring that they remain irrecoverable is known as Secure 
Data deletion. It has currently become one of the most pressing concerns in the 
cybersecurity, regulatory compliance, and digital forensics domains. While there exists 
physical means of making data irrecoverable in the physical world where confidential 
documents are shredded, and storage media containing classified material are physically 
destroyed, in the digital domain this is significantly complex. The standard delete 
operation on modern systems only implements file unlinking, which only changes the file 
system metadata while full contents will still remain intact in the system but users 
assume that all the content has been permanently erased. \cite{reardon2013sok}

The criticality of this problem has been addressed by two intersecting forces. First, 
regulatory frameworks including GDPR right to erasure, NIST SP 800-88 Rev. 1, and IEEE 
2883-2022 have mandated demonstrative and verifiable data removal evidences at defined 
sanitization levels; Clear, Purge and Destroy, with specific assurance against data recovery. 
Next, the mainstream transition from the usage of Hard Disk Drives (HDDs) to Solid-State 
Drives (SSDs) has led to a fundamental alteration of the storage landscape resulting in 
decades of secure deletion practices being inadequate.\cite{aronsson2025secure} The 
problem of secure deletion not only occurs in the technical domain but also spans into
the legal and ethical domains. On the legal front, regulations such as the GDPR, the 
California Consumer Privacy Act (CCPA), and emerging AI governance frameworks impose 
obligations for data erasure that existing technical mechanisms struggle to satisfy.

Taking into account the criticality of the problem, a Study conducted by two MIT graduate 
students in 2003 by purchasing 158 used hard drives from eBay found that 117 (74\%) drives 
contained recoverable data, 28 (17\%) had fully installed operating systems with user data 
requiring no special effort to recover, only 12 drives (8\%) had been properly sanitized. 
Recovered data included medical records, love letters, corporate financial records, and 
over 5,000 credit card numbers \cite{MITstudy2003}. A near-identical experiment was conducted
by Blancco Technology Group in 2018 with 159 used drives across the U.S., U.K., Germany, and 
Finland. Every seller insisted the drives had been "wiped." Researchers recovered passport 
scans, spreadsheets full of customer data, and intimate photographs. \cite{NetGuardia2026}.

As per the 2024 Verizon Data Breach Investigations Report, 68\% of data breaches involve 
non-malicious human error—including failure to properly delete sensitive information from 
decommissioned devices \cite{Verizon2024DBIR}. Separately, Blancco Technology Group's annual 
data sanitization report found that 42\% of used drives sold on secondary markets contain 
recoverable personally identifiable information (PII), with 15\% containing corporate financial 
data \cite{Blancco2026}. The public sector contributes some of the earlist cases, In 2002, the 
Pennsylvania Department of Labor and Industry sold computers to local resellers without 
sanitizing the drives. The machines contained "thousands of files of information about state 
employees" that had never been removed \cite{MITstudy2003}. More recently, in 2026, Cornwall 
Council in the UK suffered a redaction failure when it sent PDF documents that appeared redacted 
but revealed the full personal details of 10 complainants—names, home addresses, emails, and 
phone numbers—when opened. The failure mechanism was using visual overlays, a redaction approach
that does not securely delete the underlying data instead of applying permanent redaction with 
sanitization \cite{ComplianceHub2026}.

Within the financial Sector, Between 2015 and 2022, Morgan Stanley Smith Barney (MSSB) committed 
a series of secure deletion failures that became one of the most expensive examples of improper 
data disposal on record. In 2015, MSSB decommissioned thousands of hard drives and servers by hiring
a company which had no prior experience or expertise in proper data destruction. These hard drives
and server were later resold on an internet auction site while containing the PII of 15 million 
customers. \cite{Goodin2022}. In a separate incident, 42 servers went missing during a hardware 
refresh program. All potentially contained unencrypted customer data. The servers had encryption 
capability, but the firm never activated the encryption software. Recovered drives contained at 
least 140,000 pieces of customer PII on 13 out of 14 hard drives examined forensically. The SEC 
stated the "vast majority" of devices were never recovered \cite{SEC2022}. The company ended up 
paying over \$155 million in fines and settlements related to these incidents \cite{CybersecurityDive2022}. 

The common thread across every case is the same: standard deletion does not remove data. Operating 
systems typically only remove file system pointers, leaving the underlying data intact on the 
storage medium. On SSDs, the problem compounds because the Flash Translation Layer performs 
out-of-place updates, wear leveling, and garbage collection that can silently preserve copies 
even after overwrite attempts. Whether the mechanism is a visual redaction overlay, a missing 
encryption activation, an unverified third-party promise, or a simple failure to wipe before 
resale, the result is identical: data thought to be gone remains recoverable.

The literature review below surveys state-of-the-art techniques in secure deletion, 
identifies methodological gaps, and motivates the need for a systematic framework that can 
reliably assess deletion effectiveness across diverse modern storage environments. The 
remainder of this thesis is organized as follows. The next section reviews the background of NAND 
flash storage architecture and the secure deletion problem in detail. Chapter 2 presents a 
systematic review of existing secure deletion techniques and their empirical evaluations along 
with a review on the Legal and Ethical Aspects. Chapter 3 describes the proposed experimental 
framework, including its metric system, test methodology, and verification protocols. Chapter 4
 presents the results of applying the framework across a diverse set of storage devices and 
 deletion techniques. Chapter 5 discusses the implications of the findings and Chapter 6 
 concludes with directions for future work.

\section {NAND flash storage architecture}
NAND flash memory which was introduced by Toshiba in 1988, is a type of non-volatile 
semiconductor memory that has become the dominant storage technology for solid-state drives 
(SSDs), USB flash drives, and mobile devices \cite{aronsson2025secure}. Unlike NOR flash, 
which supports random access and high read performance at smaller capacities, NAND flash 
offers significantly higher storage density and faster write/erase operations at a lower 
cost, making it ideal for bulk data storage \cite{homaidi2009data}

\subsection{Cell Structure and Memory Organization}
At the physical level, each NAND flash memory cell is a floating-gate transistor. Data is 
stored by injecting electrical charge into the floating gate, which is electrically isolated 
and can retain charge without power. A charged floating gate (programmed state) represents a 
logical "0," while an uncharged gate (erased state) represents a logical "1" \cite{homaidi2009data}.
The amount of charge determines the cell's threshold voltage (Vt), which is read and interpreted as 
one or more bits depending on the cell type:Single-level cell (SLC) flash stores one bit per cell, 
whereas multi-level cell (MLC) stores two bits, triple-level cell (TLC) stores three bits, and 
quad-level cell (QLC) stores four bits. Each increase in bit density reduces cost per gigabyte at 
the expense of reduced endurance and lower performance \cite{aronsson2025secure}.

Memory cells are arranged hierarchically. Cells connected to the same word line form a page, 
which is the smallest unit for read and program (write) operations. Pages are grouped into 
blocks, the smallest unit for erase operations. Blocks are organized into planes, multiple 
planes reside on a single die, and multiple dies are packaged into a NAND chip. Typical page 
sizes range from 2 KB to 16 KB, and blocks typically contain 64 to 256 pages\cite{aronsson2025secure}.
Each page also includes a spare or out-of-band (OOB) area used to store error-correcting codes 
(ECC) and metadata \cite{duplicateschen2022}. During a normal program operation, charge is 
injected into the floating gate using Fowler-Nordheim tunneling, shifting the Vt to a target 
distribution window. A read operation applies a reference voltage and senses whether current 
flows, revealing the cell's programmed state. The erase operation removes charge by applying a 
high voltage across the transistor, returning the cell to its erased (lowest Vt) state.

\subsection{Basic Operations and Timing}
The three fundamental operations on NAND flash are read, program, and erase. Read and program 
operate at page granularity, while erase operates at block granularity. Their latencies differ 
by orders of magnitude: a page read takes approximately 20 to 60 $\mu$s, a page program takes 200 
to 500 $\mu$s, and a block erase takes 1.5 to 5 ms \cite{aronsson2025secure,duplicateschen2022}. An 
erase-before-write architecture means that updating even a single byte requires eventually 
erasing an entire block, which the Flash Translation Layer (FTL) works around using out-of-place 
updates. The FTL abstracts these physical details behind a standard block-device interface, 
mapping logical block addresses (LBAs) to physical NAND pages \cite{onarlioglu2018_3}. This 
abstraction allows traditional file systems (ext4, NTFS, HFS+) to remain compatible, but it also 
hides the true physical state of the data from the host \cite{aronsson2025secure}.


\subsection{Write-Once Property and Out-of-Place Updates}
NAND flash has a fundamental constraint: a page cannot be overwritten in place. To write new data to a 
page, the entire encompassing block must first be erased, which resets all cells in that block to the 
"1" (erased) state \cite{aronsson2025secure} Because erasing a block is an expensive operation that 
destroys all data within it, SSDs do not perform in-place updates. Instead, they use out-of-place 
updates: modified data is written to an available free page elsewhere, and the old page is marked 
as invalid \cite{zhang2018ensuring}. Pages containing current data are termed valid pages, while 
those holding superseded data are invalid pages. This write-once property is the root cause of many 
of the secure deletion challenges discussed in this review.

\subsection{Flash Translation Layer}
To maintain compatibility with traditional block-based file systems (e.g., ext4, NTFS), SSDs 
incorporate a Flash Translation Layer (FTL) that runs on the device controller \cite{aronsson2025secure}.
The FTL maps logical block addresses (LBAs) issued by the operating system to physical flash 
addresses. This abstraction hides the complexity of NAND flash from the host and enables the 
SSD to manage internal operations transparently \cite{duplicateschen2022}.

The FTL implements several critical management functions:

\textbf{Wear leveling} distributes program/erase (P/E) cycles evenly across all blocks to extend 
device lifespan. Each NAND flash block has a finite endurance typically 100,000 cycles for SLC 
and 10,000 to 100,000 for MLC/TLC after which the block becomes unreliable. Dynamic wear leveling 
distributes writes among free blocks, while static wear leveling periodically relocates infrequently 
accessed (cold) data to prevent some blocks from remaining idle while others wear out \cite{homaidi2009data}.

\textbf{Garbage collection} reclaims space occupied by invalid pages. When the number of free pages falls 
below a threshold, the FTL selects a victim block with a high proportion of invalid pages, copies any 
remaining valid pages to a free block, and then erases the victim block \cite{aronsson2025secure,duplicateschen2022}.
This process introduces write amplification, where a single host write triggers multiple internal writes 
due to data migration.

\textbf{Bad block management} tracks blocks that have become defective through wear or manufacturing defects. 
These blocks are recorded in a bad block table and excluded from further use, with any valid data relocated 
to healthy blocks \cite{duplicateschen2022,homaidi2009data}.

\textbf{Over-provisioning} reserves a portion of the total flash capacity as spare area to improve performance 
and endurance. This over-provisioned space is used for wear leveling, garbage collection, and bad block management. 
\cite{aronsson2025secure,ahn2021forensics}

\textbf{Error correction} is essential because NAND flash cells are susceptible to read disturb, program disturb, 
charge leakage, and retention errors. The FTL computes ECC data stored in each page's OOB area to detect and 
correct bit errors\cite{zhang2018ensuring}.

An SSD packages multiple NAND chips with a controller that manages the FTL, host interface, and error correction. 
The storage architecture spans several layers: the host interface (e.g., SATA or NVMe) communicates with the FTL, 
which translates LBAs to physical addresses and coordinates access across multiple channels, packages, and dies 
in parallel to maximize throughput. To overcome the density scaling limits of planar (2D) NAND, manufacturers 
have adopted 3D NAND technology, which stacks memory cells vertically in layers \cite{aronsson2025secure}. 
3D NAND achieves higher storage density without shrinking individual cell geometries, which improves endurance 
and reduces interference between adjacent cells. However, the vertical architecture introduces more complex 
disturb characteristics and new challenges for secure deletion, as erasing data in one layer can disturb 
adjacent layers \cite{ahn2021forensics}.

\subsection{Scrubbing: Overwriting at the Page Level}
Since page-level erase is physically impossible, scrubbing was proposed by Wei et al. in 2011 as a 
way to destroy data on individual programmed pages without block erasure. Scrubbing works by 
re-programming all remaining 1 bits on a page to 0, effectively overwriting the entire page with 
zeros. This is possible because NAND flash does allow programming individual bits from 1 to 0; the 
reverse transition (0 to 1) is what requires a block erase. The mechanism uses incremental step 
pulse programming (ISPP), which iteratively applies programming pulses and verifies the cell Vt 
until it reaches the target state \cite{aronsson2025secure,3dnandwang2018}.

Scrubbing comes with a serious reliability cost called cell disturbance. The act of reprogramming 
one cell unintentionally shifts the Vt of neighboring cells due to parasitic coupling through the 
word lines and bit lines. In 2D NAND, programming a page disturbs the two immediately adjacent 
pages. In 3D NAND, where flash cells are stacked vertically, as many as four adjacent pages can be 
disturbed, and the problem grows more severe with each additional layer. The increased bit error 
rate (BER) from scrubbing disturbance may trigger error correction failures, requiring the data on 
affected adjacent pages to be migrated before scrubbing can safely proceed.

The concept of a scrub budget was introduced to limit how many times a given block can be scrubbed 
before exhibiting significant data error risk. Once the budget is exceeded, secure deletion for 
that block must fall back to a full block erase or garbage collection. Determining the correct 
budget requires experimental BER analysis and is pattern-dependent and wear-cycle-dependent, making 
it difficult to set a universal safe threshold\cite{antiforensics2016,zhang2018ensuring}.

A full scrubbed page is uniformly zeros, which itself signals to a forensic examiner that a 
deletion occurred at that location. Partial scrubbing addresses this by modifying only a subset of 
bits in the page or block, avoiding the tell-tale all-zero pattern. This is a form of anti-forensic 
deletion that blends sanitized data with the appearance of normal programmed content.

More critically, Hasan and Ray demonstrated in 2020 that scrubbing is partially reversible\cite{Fukami2023}. By 
performing analog reads that measure the precise threshold voltage distribution (rather than a 
single binary interpretation), an attacker can distinguish between a cell that was programmed to a 
1 and then scrubbed to 0, and a cell that was originally programmed to 0. The residual Vt 
difference, though small, is detectable with precision sense amplifiers. Analog sanitization was 
proposed as a remediation: instead of a single binary scrubbing pass, the page is reprogrammed to a 
Vt distribution that is statistically indistinguishable from a native 0-programmed page\cite{ahn2022security}.

\subsection{Partial Overwriting and Down-Bit Programming}
Multiple finer-grained overwriting schemes operate at sub-page granularity by exploiting the fact 
that MLC and TLC cells have multiple threshold voltage states. Partial overwriting programs random 
data into only some of the Vt states of a multi-bit cell, destroying the original data without fully 
programming the cell to the highest state. Down-bit programming takes this further by programming a 
page originally stored in a 3-bit-per-cell (TLC) mode into a 1-bit-per-cell (SLC) mode, effectively 
compacting and destroying the original data in a single operation. The deletion pulse application 
scheme applies a series of erase-like pulses to a programmed page, shifting its Vt distribution 
without performing a full block erase. These schemes cause lower program disturbance than 
conventional scrubbing because they do not push cells to the highest Vt state, but they also provide 
less security assurance since some Vt residue may remain \cite{ahn2019schemes,ahn2022security}. 

\subsection{The Communication Path}
The operating system does not have direct knowledge of the physical flash layout. After a file is 
deleted, the file system deallocates the logical blocks by removing their entries from the file 
system index (inode table, MFT, etc.), but the actual data remains intact on the NAND cells. The host 
cannot simply overwrite the physical location because the FTL may have remapped those logical 
addresses to any physical page on the drive. There are three main channels through which the OS 
conveys block-deallocation information to the storage device

\textbf{TRIM} : The TRIM command is an ATA command that allows the file system to inform the SSD 
controller about which LBAs no longer store valid data. When a file is deleted, the OS issues a TRIM 
command for the LBAs covered by that file, and the controller adds those blocks to an internal queue 
of blocks to be blanked. The TRIM command was not designed for secure deletion but rather as an 
efficiency optimization for flash: without it, the FTL cannot distinguish logically discarded blocks 
from valid data and must assume that the only deletable block on the physical medium is the one being 
overwritten, eventually causing a thrashing effect. TRIM operates through the ATA DATA SET MANAGEMENT 
command and is supported by most modern operating systems (Windows 7+, Linux 2.6.28+, macOS) and file 
systems (NTFS, ext4, HFS+). Importantly, the TRIM command provides no false negatives: it indicates 
every time a block is discarded. However, it is an asynchronous notification: the controller is free 
to schedule the actual physical erasure at its own discretion, and the OS receives no confirmation of 
when or whether the erasure happened. TRIM is a suitable command for deleting specific files, but the 
FTL only uses the information to mark address mappings as invalid \cite{reardon2013sok,trimwearlevelling2019}.

\textbf{TrueErase} : Diesburg et al. proposed TrueErase as a refinement that provides similar 
information to TRIM but only for blocks belonging to files specifically marked as sensitive. 
TrueErase adds a new communication channel between the file system and the device driver that 
forwards deallocation information selectively. Device drivers are modified to implement immediate 
secure deletion when provided with this notification, avoiding the performance cost of blanket TRIM 
for all deleted blocks \cite{reardon2013sok}.

\textbf{Out-of-place updates as implicit notification} : When the host writes new data to an LBA that 
already has a physical mapping, the FTL performs an out-of-place update: it allocates a new physical 
page, writes the new data there, and marks the old physical page as invalid. This gives the FTL 
indirect evidence that the old data at that LBA is no longer needed. However, this mechanism only 
works for blocks that are explicitly overwritten, not for blocks that are simply deallocated by file 
deletion or truncation \cite{antiforensics2016}.

The FTL performs the address mapping function that translates the logical addresses used by the host 
into the physical addresses on the NAND flash memory. When the host issues a deletion or overwrite 
request, the FTL invalidates the address mapping information linking the logical address to the 
previous physical address. For deletion, the FTL receives the file's LBA range through the TRIM 
command and marks the corresponding address mappings as invalid \cite{antiforensics2016}.


Crucially, the invalidation step is purely a metadata operation. The FTL terminates the host's 
request immediately after updating the mapping table, and the actual data on the NAND cells is left 
entirely untouched. The flash controller does not erase, overwrite, or otherwise modify the physical 
pages at this point. The data persists in an undamaged state on the NAND cells, fully recoverable 
through low-level access techniques. The only change is that the FTL will never again return that 
physical page in response to a read of that LBA. This design is intentional: erasing a block is 
time-consuming (on the order of milliseconds per block), and performing it eagerly on every 
deallocation would destroy performance\cite{duplicateschen2022}. 


\section{Implications for Secure Deletion}

The architectural features described above create fundamental obstacles to secure deletion. Out-of-place updates 
cause multiple copies of data to persist across the device as stale pages are left behind. Garbage collection 
and wear leveling can relocate data transparently, making it impossible for host-level software to predict where 
data physically resides. Over-provisioned areas are inaccessible to the operating system, yet may retain residual 
copies of deleted data. These characteristics mean that a conventional delete command at the file-system level 
does not remove the underlying data from the physical medium, and techniques that are effective on HDDs such as 
overwriting are ineffective on NAND flash \cite{aronsson2025secure}.

\subsection*{Garbage Collection and the Zombie Time}
Garbage collection is the FTL process that physically reclaims NAND blocks containing invalid pages. 
Its canonical procedure has three steps: selecting a victim block with the largest number of 
invalid pages, copying any remaining valid pages from the victim block to a free block, and then 
erasing the victim block \cite{duplicateschen2022}. The erase operation makes the block available 
for future writes. Free blocks are managed from a pool, and modern SSDs are typically 
over-provisioned with 25\% or more spare NAND capacity to absorb write pressure and give garbage 
collection time to keep up \cite{onarlioglu2018_3}. 

The zombie time is defined as the interval between the point when the FTL invalidates the address 
mapping of a file's pages and the point when garbage collection physically erases those pages from 
the NAND cells \cite{antiforensics2016}. During this window, the data remains fully intact and 
recoverable, even though the file system and the host consider it deleted. The zombie time has no 
fixed duration: it depends on the storage usage pattern, the prevalence of free blocks in the 
over-provisioning pool, the FTL's victim selection policy, and whether TRIM commands were issued for 
the deallocated blocks.

The fundamental cause of the zombie time is that the FTL is designed to minimize garbage collection 
frequency, since garbage collection degrades I/O response by introducing read-copy-erase cycles 
that add latency. The FTL only triggers garbage collection under two conditions. In the foreground 
(normal garbage collection), garbage collection is triggered when the free-block pool drops below a threshold, 
meaning the device has no clean blocks available to service incoming writes. n the background, some 
FTL implementations opportunistically erase stale blocks during idle time. The foreground garbage collection event 
is driven entirely by write pressure: if the host is mostly reading and the free-block pool is 
adequate, garbage collection may not occur for long periods. There is a compounding factor: even 
when normal garbage collection is triggered, the zombie time for a specific file depends on whether the blocks 
containing its pages are selected as victim blocks. The FTL selects the victim with the highest 
count of invalid pages \cite{antiforensics2016}, so a block containing only a few invalid pages is 
unlikely to be chosen until many pages in that block have accumulated. If a block contains mostly 
valid pages plus one or two invalidated pages from a deleted file, that block will sit untouched 
indefinitely while the deleted data remains recoverable. 

The Zombie-Kill (ZK) TRIM scheme, proposed by Kwak et al., introduces a compulsory garbage 
collection mechanism that directly limits the zombie time, it operates in three modes that trade 
deletion speed against I/O performance \cite{antiforensics2016}. \textit{Normal mode} defers all physical 
erasure to the FTL's standard garbage collection policy (normal garbage collection). The TRIM handling process is 
the same as standard file deletion: the address mapping is invalidated, and the target pages are 
erased only when the free-block pool runs low and the blocks happen to be selected as victims. This 
mode has the longest zombie time, and in workloads with sufficient free blocks, the zombie time 
becomes effectively unbounded: deleted data may persist indefinitely because garbage collection is 
never triggered. \textit{Immediate mode} erases the data using the compulsory garbage collection upon receiving the 
TRIM command. As soon as the TRIM command arrives, the FTL invalidates the address mappings and 
then immediately performs compulsory garbage collection on the blocks containing the target pages. This mode limits 
the zombie time to the execution time of the compulsory garbage collection itself, which is determined by the 
number of blocks containing target pages and the number of valid pages that must be copied out of 
each victim block. In the experiments by Kwak et al., immediate mode achieved the smallest zombie 
time in all cases, but the I/O response time was 68 times larger on average than normal mode because 
the compulsory garbage collection introduces processing delays even when free blocks are abundant. \textit{Suspension 
mode} balances the trade-off. After receiving the TRIM command, the FTL saves the invalidation time 
of each deleted file's pages to a TRIM command table. The FTL then monitors the suspension time of 
each file by comparing the current time against the invalidation time. Until the suspension time 
exceeds a configurable maximum, target pages can only be erased by the normal garbage collection 
policy. If target pages remain unerased after the maximum suspension time, the blocks containing 
them are erased by the compulsory garbage collection. This mode limits the zombie time to the sum of the maximum 
suspension time and the compulsory garbage collection execution time, while producing a smaller I/O penalty than 
immediate mode. Fewer compulsory garbage collection events occur because target pages accumulate during the 
suspension window, and blocks with a higher proportion of invalid pages are selected, making each 
garbage collection event more efficient \cite{antiforensics2016}.

The zombie time is workload-dependent with significant forensic implications. In write-heavy 
workloads, the free-block pool drains quickly and normal garbage collection is triggered frequently, reducing the 
zombie time. In read-heavy workloads, the free-block pool remains ample, and normal garbage collection may never 
activate the blocks containing deleted pages. Under these conditions, the zombie time can be 
effectively infinite: the deleted data persists until the device is eventually subjected to a write-
intensive workload that forces garbage collection. Compounding this, the distribution of invalid 
pages across blocks affects whether a given deleted file's pages are reclaimed. The FTL selects the 
victim with the largest number of invalid pages, so blocks where the deleted file's pages are 
scattered across many blocks (each containing few invalid pages) experience longer zombie times 
than blocks with concentrated invalid pages.

\subsection*{TRIM Implementation Variance Across Drives}: 
Black-box testing has revealed significant variation in how different SSDs implement TRIM and 
garbage collection. On the Corsair F60, TRIM triggered within 1-10 seconds of a deletion. On other 
drives, TRIM behavior depended on the percentage of allocated space. Some drives did not erase all 
TRIM-marked blocks but rather a subset based on internal drive policies, meaning that even after 
TRIM and GC, some deleted blocks remained unerased. Importantly, garbage collection can operate 
independently of the host: it can be triggered whenever the disk is powered on, even if no write 
commands are issued, rendering write blockers and other forensic precautions useless. This means 
the forensic examiner cannot assume that data persists in a known state across reboots. Most modern 
operating systems and their default file systems have built-in support for the ATA TRIM command, 
but the implementation of TRIM varies from controllers in SSDs to kernels in operating systems. \cite{Bonetti2014,trimwearlevelling2019}

\subsection*{FTL Internals and its Forensic Impact}
Wear leveling distributes program-erase (P/E) cycles evenly across all blocks to prevent premature 
wear-out of frequently written blocks. Dynamic wear leveling operates on newly allocated blocks, 
while static wear leveling also moves cold (infrequently rewritten) data to ensure even wear on 
otherwise idle blocks. Every data migration performed by wear leveling creates opportunities for 
data duplication: a page may be copied to a new physical location while the old location is not 
immediately erased, leaving recoverable copies that neither the file system nor the OS is aware of. 
The write amplification factor (WAF) measures the ratio of physical writes to logical writes. A WAF 
of 3 means that every logical write from the host causes 3 physical writes due to garbage 
collection and wear leveling overhead. Each extra write represents a potential duplicate of the 
original data, and each duplicate must be separately sanitized for secure deletion to be effective.
\cite{duplicateschen2022,onarlioglu2018_3}

The copy-back program operation is an internal NAND flash operation that moves data from one page to
another without transferring data to the controller. It is used by FTLs during garbage collection 
and wear leveling to reduce latency. Ahn and Lee revealed that the copy-back program operation 
leaves unmanaged original data in the source page, which becomes a major forensic target. The 
source page may contain the original data even after the copy-back completes, and because the 
copy-back is a low-level operation invisible to the host, no higher-layer deletion mechanism can 
target it. Secure copy-back programming adds an enhanced secure program operation after the 
conventional copy-back. This can take two forms: programming random data into the source page 
(destroying the stored data) or applying a series of deletion pulses to the source page to shift 
its Vt distribution. To minimize disturbance to adjacent victim cells, the controller can first 
read victim cell group information from the adjacent word line and selectively inhibit programming 
on cells that would cause errors. \cite{ahn2021forensics}

FTL operations systematically create multiple copies of data that persist after the primary copy is 
deleted. Garbage collection copies valid pages from victim blocks to free blocks before erasing, 
leaving the original data in unerased cells that may remain in the free block pool. Wear leveling 
similarly migrates data, and the old physical location is not always erased immediately. Bad block 
management copies data away from failing blocks, creating yet another duplicate. The RedFlash 
scheme addresses this by maintaining a duplicate chain in the out-of-band (OOB) area of each flash 
page. The OOB area is typically reserved for ECC parity and bad block markers, the RedFlash was 
repurposed to store metadata linking each page to all other pages that contain copies of the same 
logical data. When a deletion command is issued, the OOB chain enables backtracking so that every 
duplicate is located and erased, not just the primary copy \cite{duplicateschen2022}.

Encryption-based deletion alone does not eliminate structural artifacts. Chen and Chen showed that 
the physical location of data on flash storage can itself leak information even when the data 
content is encrypted or destroyed. For example, knowing that certain blocks are allocated to 
specific applications or user profiles reveals usage patterns. When a block is sanitized while 
surrounding blocks remain intact, the spatial pattern of erasures implicitly confirms which data 
the user considered sensitive. Sanitizing data content without addressing structural context leaves 
a recoverable forensic profile \cite{aronsson2025secure}.

\subsection{Verification of Secure Deletion}

A persistent gap in secure deletion is the inability to independently confirm that data has been 
physically destroyed. Ahn and Lee proposed an in-band verification technique based on cell count 
information. After a secure deletion operation, the NAND cells in the sanitized area exhibit a 
characteristic Vt distribution that differs from normal programmed data. Under normal program 
conditions, the number of cells in each Vt state is approximately uniformly distributed. After a 
scrubbing or block erase, the distribution shifts to a non-uniform pattern. A simple read operation 
can count cells per distribution and compare against a reference threshold, determining pass or 
fail without comparing personal data directly. This verification can be performed in real time 
immediately after secure deletion and requires no structural changes to existing NAND flash memory.
For higher assurance, the verification scheme can use XOR operations between the data sensed from 
the sanitized area and the original personal data, with pass/fail determination accounting for the 
ECC correction capability. If the residual data is within the correctable bit error range, deletion 
is considered successful. \cite{ahn2021forensics}

Modern SSDs can vary the size of the over-provisioning (OP) area in real time to optimize 
performance, with OP areas reaching up to 50\% of total NAND capacity. Changing the OP area size 
generates an invalid data region: when OP shrinks, blocks that were previously reserved become 
accessible as user space, and the data they contain has never been sanitized. A forensic attacker 
can exploit this by manipulating the OP area allocation to expose blocks that still hold residual 
data. Even if the user data is encrypted, an attacker with sufficient resources could read the 
encrypted data from the invalidated area and attempt decryption offline. This makes deletion 
verification for dynamically allocated OP areas an essential but largely unaddressed requirement. 
Some proposed solutions include pseudo-erase operations that combine partial programming with 
lightweight verification, trading full sanitization for real-time performance \cite{ahn2021forensics}.

Achieving secure deletion in modern systems is fundamentally challenging because storage architectures are organized
into multiple layers application, file system, block device, and physical medium each of which can retain residual 
copies of data \cite{duplicateschen2022}. Performing deletion at the application layer, for example, only removes the 
logical reference; the underlying data remains in the file system and may persist on the physical media indefinitely \cite{zhang2018ensuring}.


\section{Native Drive-Level Sanitization}

Native drive-level sanitization refers to built-in hardware/firmware commands defined by storage 
interface standards (ATA, NVMe, SCSI) that instruct a drive's controller to erase all user data on 
the physical medium. Unlike software-based overwriting tools that operate through the file system, 
these commands execute inside the drive's controller and can potentially reach areas the OS cannot 
access — including over-provisioned zones and retired blocks \cite{antiforensics2016}. 

The three major standrds; NIST SP 800-88 Rev. 1 (2014), IEEE 2883-2022, and NATO AEDP-3 use a three 
tier classification of sanitization methods: Clear, Purge, and Destruct. \cite{aronsson2025secure}.

\begin{table}[H]
    \centering
    \caption{Comparison of NIST SP 800-88 Sanitization Levels}
    \label{tab:sanitization-levels}
    \begin{tabularx}{\textwidth}{lXXl}
        \toprule
        \textbf{Level} & \textbf{Description} & \textbf{Effectiveness} & \textbf{Reusability} \\
        \midrule
        Clear & Logical sanitization that removes user-accessible data only. Example: ATA Secure Erase, NVMe Format without CRYPTO flag. & Low — residual data may remain in hidden areas & Storage can be reused \\
        \midrule
        Purge & Advanced sanitization making recovery infeasible even under lab conditions. Example: ATA Enhanced Secure Erase, NVMe Secure Erase (CRYPTO flag enabled), Cryptographic Erase. & High & Storage can be reused \\
        \midrule
        Destruct & Physical destruction of the media. & Complete & Device cannot be reused \\
        \bottomrule
    \end{tabularx}
\end{table}

\subsection{Key Native Commands by Interface}

\textbf{ATA/ATAPI} 
\begin{itemize}
    \item \texttt{SECURITY ERASE PREPARE} : prerequisite that unlocks the drive for erasure.
    \item \texttt{SECURITY ERASE UNIT} : overwrites every NAND cell with zeros or ones. Operates at the factory-reset level and covers reserved areas (over-provision and retired blocks) that the host cannot normally address.
    \item \texttt{ENHANCED SECURITY ERASE} : vendor-defined extension that may erase more thoroughly (e.g., SanDisk calls this "sanitize").
    \item \texttt{SANITIZE DEVICE} : a newer, more flexible command supporting three modes: crypto scramble, overwrite, and block erase.
    \item \texttt{DATA SET MANAGEMENT} : not a sanitization command, but marks logical blocks as free so the controller's garbage collection can reclaim them over time.
\end{itemize}

\textbf{NVMe}
\begin{itemize}
    \item \texttt{Format NVM} : reformats the namespace. Without the CRYPTO flag it is a Clear-level operation; with the CRYPTO flag enabled it reaches Purge level.
    \item \texttt{Sanitize} : supports Block Erase, Overwrite, and Crypto Erase modes. The \texttt{nvme sanitize-log} command reports progress and completion status.
\end{itemize}

\textbf{SCSI}
\begin{itemize}
    \item \texttt{SANITIZE} : functionally similar to the ATA SANITIZE DEVICE command, supporting block erase, overwrite, and cryptographic erase options.
\end{itemize}

As stated above there are two principle purge mechanisms: block erase and cryptographic erase. 

\textbf{Block Erase} : The controller issues the NAND flash erase command for every block on the 
device, forcing each block to its unprogrammed (erased) state. This is inherently slow — a 900 GB 
NVMe drive can take significant time — and the time estimates reported by tools are often 
inaccurate. Critically, block erase behavior is governed by usage-aware controller logic that may 
compromise the guarantee of full sanitization; studies found that the controller sometimes returns 
all-zero data during verification without having physically erased every cell, suggesting internal 
masking rather than true physical deletion. \cite{aronsson2025secure}

\textbf{Cryptographic Erase} : On self-encrypting drives (those supporting TCG Opal SSC or similar), 
all user data is encrypted by a media encryption key stored in the controller. Cryptographic erase 
works by scrambling or deleting that key, making the ciphertext on the medium permanently 
undecryptable. This is nearly instantaneous and does not wear the NAND cells. The flip side: 
security rests entirely on the vendor's key management. If the key is not fully destroyed, if a 
copy persists in a controller cache, or if there is a backdoor, the protection fails. \cite{aronsson2025secure}

As stated above,NVMe's Sanitize command supports Cryptographic Erase as one of its three modes. Cryptographic 
erasure renders data irrecoverable by destroying the media encryption key (MEK) rather than overwriting 
physical cells and n a self-encrypting drive, all user data is encrypted by a MEK that never leaves the 
controller's secure boundary. An NVMe self-encrypting drive complying with TCG Opal 2.0 maintains 
multiple keys at different privilege levels and the MEK is the root of every data confidentiality guarantee. 
All user data written to the NAND medium is ciphertext under this single key, encrypted and decrypted 
transparently by the hardware encryption engine in the drive controller \cite{aronsson2025secure}.

\begin{table}[htbp]
    \centering
    \caption{The Key Hierarchy Within an NVMe SED}
    \label{tab:opal-keys}
    \begin{tabularx}{\textwidth}{lXX}
        \toprule
        \textbf{Key} & \textbf{Role} & \textbf{Storage Location} \\
        \midrule
        Authentication Key (AK) & Proves the user/host is authorized to issue commands & Derived from user password, never stored in plaintext \\
        \midrule
        Band Master Key (BMTK) & Protects individual locking ranges (bands) on the drive & Encrypted under the AK, stored on the NAND medium \\
        \midrule
        Media Encryption Key (MEK) & The actual AES cipher key that encrypts/decrypts all user data on the NAND flash & Inside the controller's dedicated key memory — never leaves the controller die \\
        \bottomrule
    \end{tabularx}
\end{table}

\subsection{What is MEK Rotation}

MEK rotation is the act of replacing the current Media Encryption Key with a new one and securely 
discarding the old key. Within the TCG Opal framework, this is accomplished through the RevertSP and 
Revert methods, which are the Opal equivalents of the NVMe Sanitize (Crypto Scramble) command \cite{aronsson2025secure}.

The procedure, as implemented by Opal-compliant drives, works as follows:
\begin{enumerate}
    \item  The host issues a RevertSP (or Revert) command to the drive's TCG Storage Interface.

    \item The controller generates a new random AES-128 or AES-256 key.

    \item The controller overwrites its internal MEK register with the new key — this is a purely volatile operation within the controller logic.

    \item The old MEK is destroyed as the register is overwritten. Since the MEK never leaves the controller in plaintext, and the controller has no mechanism to export it, the old key becomes unrecoverable.

    \item All ciphertext on the NAND medium is now permanently undecryptable. The data has not been physically erased — every bit pattern remains exactly where it was — but it can never be recovered because the cryptographic secret that produced the plaintext is gone.

    \item Optionally, per the NATO AEDP-3 standard, a block erase can follow the cryptographic erase to physically scrub the NAND cells, adding defense in depth \cite{aronsson2025secure}.
\end{enumerate}



A full block erase on a large NVMe drive (e.g. 1 TB) can take tens of minutes because it must cycle 
through every NAND block's erase operation sequentially, with verification \cite{aronsson2025secure}
while MEK rotation completes in approximately 0.22 seconds regardless of drive capacity or how much data 
is stored

But still, these approaches has their practical limitations. Firstly, vendors can be inconsistent in 
supporting and implementing these commands, with some drives ignoring them or providing false assuarances 
of completion \cite{wang2020challenges}. Further, Native commands are all-or-nothing where they sanitize 
the entire device which makes them suitable for decommissioning but useless in routine per-file deletions.
And the controller's internal logic is not transparent, making it difficult to verify that the intended 
sanitization has actually occurred, all-zero reads may reflect the controller lying rather than physical 
erasure \cite{aronsson2025secure}. Finally, wear-leveling mechanisms spreads data across cells, garbage collection 
creates copies in unexpected locations, and over-provisioning hides blocks from the host, all of which mean that 
even a correctly executed sanitize command may leave data remnants in retired or remapped blocks \cite{aronsson2025secure}.
A verifiable alternative method SAFE (Scramble and Finally Erase) was proposed to overcome these limitations. 
SAFE encrypts all data with a key held only in the controller's volatile memory, then on receiving a 
sanitize command: (1) erases the controller key memory, (2) erases every NAND block, (3) writes a known 
verification pattern, (4) erases again, and (5) reinitializes with a new key. This provides 
controller-level verifiability but, like the other native commands, offers no selective deletion 
capability \cite{reardon2013sok}.

SAFE (Scramble and Finally Erase) essentially performs an MEK rotation combined with a follow-up block 
erase and verification pattern — a verified MEK rotation plus physical scrub \cite{zhang2018ensuring}. The 
Eraser and Holepunch schemes operate at a higher layer, using host-level master key rotation with tree-based 
or PPRF-based file key derivation to achieve the equivalent of selective "mini-MEK rotations" at file 
granularity, solving the all-or-nothing limitation of native NVMe MEK rotation \cite{aronsson2025secure,holepunch}.

% use [] to set name for ToC
%\section[Goal]{Goal} % ok with fontsize=12pt

%\section[Struttura della tesi]{Thesis Structure} % ok with fontsize=12pt
